\section{Comparison with Published SROIE Task-3 Entries}
\label{sec:competitors}

Because the advanced variant of this paper evaluates on the
\emph{canonical} ICDAR-2019 SROIE Task-3 test set (347 official
held-out receipts) every number reported here is directly comparable
to published SROIE Task-3 results from the wider literature.
Table~\ref{tab:competitors} reproduces this side-by-side: each row
either cites the original paper that reported the F1 number on the
347-image set, or — for the two ``this work'' rows — reads the live
measurement out of the run's
\texttt{combined\_metrics.json}.  The fixture file
\texttt{results/sroie\_task3\_competitors.json} is the editable
source of truth; updating a competitor's F1 only requires editing
that JSON and rerunning the paper stage.

Our headline claim --- \emph{matching} our re-implemented
DONUT~\cite{kim2022donut} on Task-3 with \VAR{param_ratio_phrase}
the parameter budget (see Sec.~\ref{sec:discussion} for the
published-vs-re-implemented framing) --- is decided by the row pair
``this work --- DONUT'' vs ``this work --- YOLO+TrOCR+Assigner''
combined with the published baselines underneath.

\begin{table}[t]
  \centering
  \caption{Global token-level F1 on the canonical SROIE Task-3 test
           set (347 images).  ``this work'' rows are sourced from
           \texttt{combined\_metrics.json}; the remainder are
           published numbers cited from each paper's results table.
           This table is emitted only under the \texttt{advanced}
           paper variant; the \texttt{basic} variant evaluates on
           the internal 500/63/63 split which is not leaderboard-
           comparable and thus omits this comparison.}
  \label{tab:competitors}
  \VAR{table_competitors}
\end{table}
